\documentclass[11pt]{article}
\usepackage[margin=1in]{geometry}
\usepackage{amsmath,amssymb,amsthm}
\usepackage{xcolor}
\usepackage[ruled,vlined]{algorithm2e}
\SetKw{KwBreak}{break}
\usepackage{enumitem}
\usepackage{hyperref}

\newcommand{\E}{\mathbb{E}}
\newcommand{\R}{\mathbb{R}}
\newcommand{\F}{\mathcal{F}}
\newcommand{\G}{\mathcal{G}}
\newcommand{\pr}{\mathbb{P}}
\newcommand{\T}{^{\!\top}}
\newtheorem{assumption}{Assumption}
\newtheorem{remark}{Remark}

\title{Deep Picard regression for Nash equilibria under asymmetric information}
\author{}
\date{}

\begin{document}
\maketitle

\section{Method}
\label{sec:method}

In this section, we present a method for stochastic differential games in which the players observe different information. The first order condition of each player expresses its control as a conditional expectation given its own observations. A Nash equilibrium is a set of controls for which these conditions hold for all players at once. We compute it by simulating the game under the current controls, approximating each conditional expectation by minimizing a mean squared error over the simulated paths, and iterating. In the remainder of this section, we describe the game, the first order condition of one player, the equilibrium as a fixed point, and the numerical scheme.

\subsection{The game}
\label{sec:game}

We work on a filtered probability space $(\Omega, \G, \mathbb{G} = (\G_t)_{t \in [0,T]}, \pr)$, where $T > 0$ is a fixed and finite horizon, and $\mathbb{G}$ is the filtration generated by a $d$-dimensional Brownian motion $W$.

There are $n$ players. Player $i$ chooses a control $u^i = (u^i_t)_{t \in [0,T]}$, $u^i_t \in \R^{d_i}$. We denote by $u = (u^1, \dots, u^n)$ the controls of all players, and by $u^{-i}$ the controls of the players other than $i$. The state of the game is denoted by $X = (X_t)_{t \in [0,T]}$, $X_t \in \R^{d_X}$, and satisfies
\begin{equation}
  \label{eq:state}
  dX_t = \mu(t, X_t, u_t) \, dt + \sigma(t, X_t) \, dW_t , \qquad X_0 = x_0 .
\end{equation}

We assume player $i$ observes a process $O^i = (O^i_t)_{t \in [0,T]}$, a function of the state, but not the state itself. We denote by $\mathbb{F}^i = (\F^i_t)_{t \in [0,T]}$ the filtration generated by $O^i$, so that $\F^i_t = \sigma\{O^i_s : 0 \le s \le t\} \subseteq \G_t$. The filtrations of different players need not be nested. A control of player $i$ is admissible if it is $\mathbb{F}^i$-progressively measurable and square integrable, and we write $\mathcal{A}^i$ for the set of admissible controls.

Player $i$ seeks to maximize
\begin{equation}
  \label{eq:objective}
  J^i[u^i, u^{-i}] = \E\!\left[ \int_0^T r^i(t, X_t, u_t) \, dt + g^i(X_T) \right]
\end{equation}
over $\mathcal{A}^i$, where $r^i$ is the running reward and $g^i$ the terminal reward.

\subsection{First order condition}
\label{sec:foc}

We fix $u^{-i}$ and assume that $J^i[\,\cdot\,, u^{-i}]$ is strictly concave on $\mathcal{A}^i$, so that the optimal control of player $i$ is characterized by a first order condition. We denote by $H^i(t, x, u, p) = r^i(t, x, u) + p\T \mu(t, x, u)$ the Hamiltonian of player $i$, and by $(P^i, Q^i)$ the adjoint processes, which solve
\begin{equation}
  \label{eq:adjoint}
  dP^i_t = - \partial_x H^i(t, X_t, u_t, P^i_t) \, dt + Q^i_t \, dW_t , \qquad P^i_T = \partial_x g^i(X_T) .
\end{equation}
Since admissible perturbations of $u^i$ are $\mathbb{F}^i$-adapted, the first order condition holds conditionally on $\F^i_t$,
\begin{equation}
  \label{eq:foc}
  \E\big[ \partial_{u^i} H^i(t, X_t, u_t, P^i_t) \,\big|\, \F^i_t \big] = 0 .
\end{equation}
The state equation \eqref{eq:state}, the adjoint equations \eqref{eq:adjoint} and the conditions \eqref{eq:foc}, for $i = 1, \dots, n$, form a forward-backward system. The forward equation is shared, and each control is determined on its own filtration. \textcolor{red}{Cite the partially observed stochastic maximum principle.}

In the examples we consider, $\mu$ is linear in $(x, u)$ and $r^i$ is quadratic in $u^i$, with $\partial_{u^i} r^i = - 2 a_i u^i_t + b^i(t, X_t, u^{-i}_t)$ and $a_i > 0$. Then $P^i$ is an explicit integral of the future rewards and \eqref{eq:foc} gives
\begin{equation}
  \label{eq:control}
  u^i_t = \frac{1}{2 a_i} \E\big[ \Lambda^i_t \,\big|\, \F^i_t \big] , \qquad
  \Lambda^i_t = b^i(t, X_t, u^{-i}_t) + \int_t^T \Phi^i(s, t)\T \partial_x r^i(s, X_s, u_s) \, ds + \Phi^i(T, t)\T \partial_x g^i(X_T) ,
\end{equation}
where $\Phi^i(s, t)$ is the deterministic sensitivity of $X_s$ to $u^i_t$. The random variable $\Lambda^i_t$ depends on the future of the game, including the future controls of player $i$. Thus, \eqref{eq:control} characterizes the optimal control but does not give it explicitly.

\subsection{Equilibrium as a fixed point}
\label{sec:fixed-point}

The controls $u^\star$ form a Nash equilibrium if $J^i[u^{i\star}, u^{-i\star}] \ge J^i[v, u^{-i\star}]$ for every player $i$ and every $v \in \mathcal{A}^i$. We denote by $\Psi$ the map
\begin{equation}
  \label{eq:psi}
  \big( \Psi(u) \big)^i_t = \frac{1}{2 a_i} \E\big[ \Lambda^i_t[u] \,\big|\, \F^i_t \big] , \qquad i = 1, \dots, n ,
\end{equation}
where $\Lambda^i_t[u]$ is \eqref{eq:control} evaluated under the controls $u$. Under the concavity assumption, $u$ is a Nash equilibrium if and only if it is a fixed point of $\Psi$. The system \eqref{eq:state}, \eqref{eq:adjoint} and \eqref{eq:foc}, for $i = 1, \dots, n$, is the corresponding forward-backward system, in which each control is determined on its own filtration. \textcolor{red}{For $n = 2$, compare with Wu and Jaimungal (2024); verify the reference.}

\subsection{Conditional expectations by regression}
\label{sec:regression}

We write $L^2(\F^i_t)$ for the space of square integrable $\F^i_t$-measurable random variables. The conditional expectation is the $L^2$-projection
\begin{equation}
  \label{eq:projection}
  \E[\Lambda^i_t \mid \F^i_t] = \operatorname*{argmin}_{\zeta \in L^2(\F^i_t)} \; \E\big[ |\zeta - \Lambda^i_t|^2 \big] .
\end{equation}
Every element of $L^2(\F^i_t)$ is a function of the observed path $(O^i_s)_{s \le t}$. On each simulated path, $\Lambda^i_t$ can be evaluated, since the whole future is known. Thus, the conditional expectation can be approximated by fitting a function of the observed path to these values by minimizing a mean squared error. We represent the function by a recurrent neural network $\mathcal{N}^i(\cdot\,; \theta^i)$, which takes the increments of $O^i$ as input and returns the control at each time step.

\subsection{Iteration}
\label{sec:iteration}

We compute the fixed point of $\Psi$ by the damped Picard iteration
\begin{equation}
  \label{eq:picard}
  u^{(k+1)} = (1 - \rho) \, u^{(k)} + \rho \, \Psi\big( u^{(k)} \big) , \qquad \rho \in (0, 1] ,
\end{equation}
in which all players are updated at once. \textcolor{red}{State conditions under which \eqref{eq:picard} converges, at least in the linear-quadratic case.}

To compute \eqref{eq:picard} numerically, fix a time grid $t_j = j \Delta t$, $j = 0, \dots, N$, with $\Delta t = T/N$. At iteration $k$, we simulate $M$ paths of $X$ under $u^{(k)}$ by the Euler--Maruyama scheme, and evaluate $\Lambda^i$ along each path by a Riemann sum. We denote the result by $\Lambda^{i,(m,k)}_j$. The parameters of $\mathcal{N}^i$ minimize
\begin{equation}
  \label{eq:loss}
  \mathcal{L}^i(\theta^i) = \frac{1}{M N} \sum_{m=1}^{M} \sum_{j=0}^{N-1} \Big| \mathcal{N}^i\big(t_j, O^{i,(m)}_{[0, t_j]}; \theta^i\big) - (1 - \rho) \, u^{i,(m,k)}_{t_j} - \frac{\rho}{2 a_i} \Lambda^{i,(m,k)}_j \Big|^2 .
\end{equation}
At each iteration, we initialize each network with the parameters of the previous iterate. After $K$ iterations, we average the controls of the last iterates.

The network of player $i$ takes only $O^i$ as input. The unobserved processes, including the other players' observations, enter only through the target $\Lambda^i$. No filter, and no estimate of the other players' estimates, is computed explicitly.

\subsection{Algorithm summary}
\label{sec:algorithm}

The computational procedure is as follows.
\begin{enumerate}[itemsep=4pt]
    \item \textit{Initialization:} Fix the time grid $t_j = j \Delta t$, $j = 0, \dots, N$. For each player $i$, create a recurrent network $\mathcal{N}^i(\cdot\,; \theta^{i,(0)})$ with randomly initialized parameters.

    \item \textit{Picard iteration loop:} For $k = 0, 1, \dots, K - 1$:
        \begin{enumerate}[label=(\alph*), itemsep=4pt]
            \item \textit{Forward simulation:} Simulate $M$ paths of the noise. Evolve the state \eqref{eq:state} by the Euler--Maruyama scheme, where at each step the control of player $i$ is the output of $\mathcal{N}^i(\cdot\,; \theta^{i,(k)})$ on the observed path $O^{i,(m)}_{[0, t_j]}$. Store the state paths, the controls $u^{i,(m,k)}_{t_j}$ and the observed paths.

            \item \textit{Targets:} For every player $i$, path $m$ and time step $t_j$, evaluate the pathwise functional $\Lambda^{i,(m,k)}_j$ of \eqref{eq:control} by a Riemann sum along the path.

            \item \textit{Regression:} For every player $i$, initialize $\mathcal{N}^i(\cdot\,; \theta^{i,(k+1)})$ with the parameters $\theta^{i,(k)}$ and train it on the inputs $\{O^{i,(m)}_{[0, t_j]}\}$ and the targets $\{(1 - \rho)\, u^{i,(m,k)}_{t_j} + \tfrac{\rho}{2 a_i} \Lambda^{i,(m,k)}_j\}$ by minimizing the mean squared error \eqref{eq:loss}.

            \item \textit{Convergence check:} Simulate validation paths with the new networks and compute the root mean square change of the controls between iterations $k$ and $k + 1$. Stop if it falls below a tolerance $\epsilon_{\mathrm{tol}}$.
        \end{enumerate}

    \item \textit{Return:} Return the networks $\mathcal{N}^i(\cdot\,; \theta^{i,(K)})$, or the average of the controls of the last iterates along the validation paths.
\end{enumerate}

\begin{algorithm}[H]
    \caption{Deep Picard regression for a Nash equilibrium under asymmetric information}
    \label{alg:main}
    \BlankLine
    \KwIn{Hyperparameters $M, N, K, \rho, S, \epsilon_{\mathrm{tol}}$; simulator of \eqref{eq:state}; observation maps $O^i$; functionals $\Lambda^i$.}
    \KwOut{Networks $\mathcal{N}^i(\cdot\,; \theta^i)$, $i = 1, \dots, n$.}
    \BlankLine
    Initialize networks $\mathcal{N}^i(\cdot\,; \theta^{i,(0)})$ with random parameters, $i = 1, \dots, n$\;
    \BlankLine
    \For{$k \leftarrow 0$ \KwTo $K - 1$}{
        Simulate $M$ paths of $X$ with $u^{i,(k)}_{t_j} = \mathcal{N}^i(t_j, O^i_{[0, t_j]}; \theta^{i,(k)})$ for all $i$\;
        \For{$i \leftarrow 1$ \KwTo $n$}{
            Compute targets for all $m, j$: $\mathcal{T}^i[j, m] \leftarrow (1 - \rho)\, u^{i,(m,k)}_{t_j} + \frac{\rho}{2 a_i} \Lambda^{i,(m,k)}_j$\;
            Initialize $\mathcal{N}^i(\cdot\,; \theta^{i,(k+1)})$ with parameters $\theta^{i,(k)}$\;
            Train $\mathcal{N}^i(\cdot\,; \theta^{i,(k+1)})$ for $S$ gradient steps on inputs $\{O^{i,(m)}_{[0, t_j]}\}$ and targets $\{\mathcal{T}^i[j, m]\}$\;
        }
        Simulate validation paths with the new networks\;
        \If{$\| u^{(k+1)} - u^{(k)} \|_{\mathrm{RMSE}} < \epsilon_{\mathrm{tol}}$}{
            \KwBreak\;
        }
    }
    \Return $\mathcal{N}^i(\cdot\,; \theta^{i,(k+1)})$, $i = 1, \dots, n$\;
\end{algorithm}

\begin{remark}
For $n = 1$, Algorithm \ref{alg:main} solves a single-agent problem under partial information without learning a filter. The latent-factor method instead learns the filter first, so that the backward process is a function of a Markov state. \textcolor{red}{State when each approach is preferable.}
\end{remark}

\end{document}
